% ================================================================
\subsection{Biểu thức chính quy (Regular expression - Regex)}

Biểu thức chính quy trên một bảng chữ cái cho trước $\Sigma$ là cách mô tả
một tập chuỗi bằng cách kết hợp các thành phần cơ bản sau:\cite{HCMUTautomataslide}

\begin{itemize}[nosep]
      \item Các chuỗi ký hiệu được tạo từ bảng chữ cái $\Sigma$.
      \item Dấu ngoặc \texttt{(}, \texttt{)} để nhóm biểu thức.
      \item Các toán tử hợp $+$, nối chuỗi $\cdot$, và bao đóng Kleene $*$.
\end{itemize}

Về mặt hình thức, một biểu thức chính quy trên bảng chữ cái $\Sigma$ được xác
định đệ quy bởi các điều kiện:

\begin{enumerate}[nosep]
      \item $\emptyset$, $\epsilon$ và mỗi ký hiệu $a \in \Sigma$ đều là biểu
            thức chính quy nguyên thủy.
      \item Nếu $r$ và $s$ là các biểu thức chính quy, thì $(r)$, $r + s$,
            $r \cdot s$ và $r^*$ cũng là các biểu thức chính quy.
      \item Một chuỗi là biểu thức chính quy khi và chỉ khi nó có thể được suy
            ra từ các biểu thức chính quy nguyên thủy bằng một số hữu hạn lần
            áp dụng các quy tắc ở trên.
\end{enumerate}

Trong ngữ cảnh project, regex được dùng như một công cụ đặc tả pattern: thay
vì liệt kê từng chuỗi có thể xuất hiện, ta mô tả cấu trúc chung của cả một lớp
chuỗi, chẳng hạn email, URL, số điện thoại hoặc ngày tháng.

\subsection{Nhận diện thực thể bằng biểu thức chính quy (Regex Detector)}
\label{sec:regex_detector}
% ================================================================

Module \texttt{regex\_detector} chịu trách nhiệm nhận diện bốn loại thực thể có
cấu trúc ký tự rõ ràng: \texttt{EMAIL}, \texttt{URL}, \texttt{PHONE} và
\texttt{DATE}. Đây là các thực thể mà mẫu chuỗi tương đối cố định, phù hợp
với phương pháp regular expression. Module trả về danh sách entity theo đúng
schema chung của project: \texttt{type}, \texttt{text}, \texttt{start},
\texttt{end}.

Về mặt thiết kế, module được tổ chức theo pipeline nội bộ: phát hiện email và
URL trước, sau đó dùng kết quả này làm exclusion spans khi phát hiện phone,
rồi dùng kết quả email, URL và phone làm exclusion spans khi phát hiện date.
Cách tiếp cận này giúp giảm thiểu false positive giữa các entity cùng module.

% ----------------------------------------------------------------
\subsubsection{Nhận diện email}
% ----------------------------------------------------------------

Biểu thức chính quy cho email tuân theo cấu trúc \texttt{local@domain.tld}:

\begin{itemize}[nosep]
      \item \textbf{Local part}: gồm chữ cái, chữ số, dấu chấm, dấu cộng, gạch
            dưới, gạch ngang.
      \item \textbf{Domain}: một hoặc nhiều nhãn phân cách bằng dấu chấm, mỗi
            nhãn chứa chữ cái hoặc chữ số.
      \item \textbf{TLD}: ít nhất 2 ký tự chữ cái.
\end{itemize}

Sau khi regex match, hệ thống loại bỏ dấu câu cuối (trailing punctuation) như
dấu chấm, dấu phẩy, dấu chấm than --- vì các ký tự này thường thuộc ngữ pháp
câu, không phải phần email.

% ----------------------------------------------------------------
\subsubsection{Nhận diện URL}
% ----------------------------------------------------------------

Pattern URL bắt đầu bằng \texttt{http://}, \texttt{https://} hoặc \texttt{www.}
và lấy tất cả ký tự liên tiếp cho đến khoảng trắng. Sau đó áp dụng hai bước
hậu xử lý:

\begin{enumerate}[nosep]
      \item Loại bỏ trailing punctuation (dấu chấm, dấu phẩy, v.v.).
      \item Cân bằng dấu ngoặc đơn: nếu URL kết thúc bằng \texttt{)} nhưng
            không chứa \texttt{(} bên trong, dấu \texttt{)} được coi là thuộc
            ngữ pháp câu và bị loại bỏ.
\end{enumerate}

% ----------------------------------------------------------------
\subsubsection{Nhận diện số điện thoại}
% ----------------------------------------------------------------

Regex cho phone bao phủ nhiều định dạng phổ biến:

\begin{itemize}[nosep]
      \item Chuỗi liên tục 8--15 chữ số (\texttt{0912345678}).
      \item Định dạng có một separator (\texttt{0475 4429797}).
      \item Mã quốc gia tùy chọn, mã vùng trong ngoặc, các nhóm số phân cách
            bằng khoảng trắng, dấu chấm hoặc gạch ngang
            (\texttt{+84 912 345 678}, \texttt{(091) 234 5678},
            \texttt{415.555.2671}).
\end{itemize}

Regex phone có phạm vi rộng nên cần hậu xử lý nghiêm ngặt để giảm false
positive:

\begin{enumerate}[nosep]
      \item \textbf{Exclusion}: bỏ qua candidate nếu span chồng lấn với entity
            email hoặc URL đã phát hiện trước đó.
      \item \textbf{Trim}: loại bỏ khoảng trắng đầu/cuối và trailing punctuation.
      \item \textbf{Kiểm tra số chữ số}: tổng số digit phải nằm trong khoảng
            $[9, 15]$. Điều kiện này loại bỏ các số ngắn (mã bưu điện, năm)
            và số quá dài.
      \item \textbf{Boundary check}: ký tự ngay trước \texttt{start} và ngay sau
            \texttt{end} không được là chữ cái, chữ số, hoặc \texttt{@}, nhằm
            tránh bắt nhầm một phần của token lớn hơn.
\end{enumerate}

% ----------------------------------------------------------------
\subsubsection{Nhận diện ngày tháng}
\label{sec:date_detection}
% ----------------------------------------------------------------

Entity \texttt{DATE} bao gồm nhiều định dạng ngày tháng phổ biến trong tiếng
Anh và tiếng Việt. Hệ thống sử dụng 11 biểu thức chính quy, chia thành bốn
nhóm.

\subsubsection*{Nhóm 1: Định dạng tiếng Việt}

\begin{itemize}[nosep]
      \item \texttt{ngày DD tháng MM năm YYYY} --- có hoặc không có prefix
            ``ngày''.
      \item \texttt{ngày D/M/YYYY} --- prefix ``ngày'' kèm dấu gạch chéo.
      \item \texttt{tháng MM năm YYYY} --- chỉ tháng/năm.
      \item \texttt{năm YYYY} --- chỉ năm, yêu cầu prefix ``năm'' để tránh
            false positive.
\end{itemize}

\subsubsection*{Nhóm 2: Định dạng tên tháng tiếng Anh}

\begin{itemize}[nosep]
      \item \texttt{January 2, 2024} / \texttt{Jan 2, 2024} --- tên tháng đầy
            đủ hoặc viết tắt, theo sau là ngày và năm.
      \item \texttt{2 January 2024} / \texttt{2 Jan 2024} --- ngày đứng trước
            tên tháng.
      \item \texttt{January 2024} / \texttt{Jan 2024} --- chỉ tháng/năm.
      \item \texttt{year YYYY} --- chỉ năm, yêu cầu prefix ``year''.
\end{itemize}

\subsubsection*{Nhóm 3: Định dạng số (ISO và phổ biến)}

\begin{itemize}[nosep]
      \item \texttt{YYYY-MM-DD} / \texttt{YYYY/MM/DD} --- ISO 8601.
      \item \texttt{DD/MM/YYYY}, \texttt{DD-MM-YYYY}, \texttt{DD.MM.YYYY} ---
            hỗ trợ cả năm 2 và 4 chữ số.
      \item \texttt{MM/YYYY} --- tháng/năm dạng số.
\end{itemize}

\subsubsection*{Nhóm 4: Năm đơn lẻ có ngữ cảnh}

Số 4 chữ số đứng một mình (ví dụ \texttt{2024}) \textbf{không} được nhận diện
là \texttt{DATE}, vì rất dễ trùng với mã số, chỉ số hoặc số lượng. Chỉ khi có
prefix rõ ràng như \texttt{năm} hoặc \texttt{year}, hệ thống mới phát hiện.

\subsubsection*{Giảm false positive cho DATE}

\begin{enumerate}[nosep]
      \item \textbf{Exclusion spans}: candidate DATE bị loại nếu chồng lấn với
            bất kỳ entity EMAIL, URL hoặc PHONE đã phát hiện trước đó. Điều
            này ngăn việc bắt nhầm segment giống ngày trong URL
            (ví dụ \texttt{https://example.com/2024/01/02}).
      \item \textbf{De-duplication nội bộ}: khi nhiều pattern cùng match
            overlapping span (ví dụ pattern \texttt{DD/MM/YYYY} và
            \texttt{MM/YYYY} cùng match trên chuỗi \texttt{01/02/2024}), hệ
            thống giữ match dài nhất.
      \item \textbf{Lookbehind/lookahead}: các pattern số dùng negative lookbehind
            \verb|(?<!\d)| và negative lookahead \verb|(?!\d)| để đảm bảo
            không bắt nhầm một phần của số dài hơn.
\end{enumerate}

% ----------------------------------------------------------------
\subsubsection{Pipeline nội bộ và xử lý overlap}
\label{sec:regex_pipeline}
% ----------------------------------------------------------------

Hàm chính \texttt{detect\_regex\_entities} gọi bốn detector theo thứ tự phụ
thuộc, truyền kết quả đã phát hiện làm exclusion cho detector tiếp theo:

\begin{enumerate}[nosep]
      \item Phát hiện \texttt{EMAIL} $\rightarrow$ danh sách $E_{email}$.
      \item Phát hiện \texttt{URL} $\rightarrow$ danh sách $E_{url}$.
      \item Phát hiện \texttt{PHONE} với exclusion = $E_{email} \cup E_{url}$
            $\rightarrow$ danh sách $E_{phone}$.
      \item Phát hiện \texttt{DATE} với exclusion = $E_{email} \cup E_{url}
                  \cup E_{phone}$ $\rightarrow$ danh sách $E_{date}$.
\end{enumerate}

Sau khi có đủ bốn danh sách, tất cả entity được gộp lại và đưa qua hàm
\texttt{\_resolve\_regex\_overlaps} để xử lý overlap cuối cùng theo priority:

\begin{center}
      \texttt{EMAIL} $>$ \texttt{URL} $>$ \texttt{PHONE} $>$ \texttt{DATE}
\end{center}

Thuật toán giải quyết overlap giống với merger tổng thể (Greedy Priority-First):
sắp xếp theo priority rồi duyệt tham lam, giữ entity nào không chồng lấn với
các entity đã giữ trước đó.

\begin{algorithm}[H]
      \caption{Pipeline nhận diện thực thể bằng regex}
      \label{alg:regex_pipeline}
      \begin{algorithmic}[1]
            \Require Văn bản đầu vào $T$
            \Ensure Danh sách thực thể $E$ không overlap, sắp theo vị trí

            \State $E_{email} \leftarrow \textsc{DetectEmail}(T)$
            \State $E_{url} \leftarrow \textsc{DetectUrl}(T)$
            \State $higher \leftarrow E_{email} \cup E_{url}$
            \State $E_{phone} \leftarrow \textsc{DetectPhone}(T,\; higher)$
            \State $E_{date} \leftarrow \textsc{DetectDate}(T,\; higher \cup E_{phone})$

            \State $all \leftarrow E_{email} \cup E_{url} \cup E_{phone} \cup E_{date}$
            \State $E \leftarrow \textsc{ResolveRegexOverlaps}(all)$
            \State \Return $E$
      \end{algorithmic}
\end{algorithm}

% ----------------------------------------------------------------
\subsubsection{Minh họa hoạt động}
% ----------------------------------------------------------------

Xét đoạn văn bản:

\begin{center}
      \textit{``Email john@a.com, born 01/02/2024, call +84 912 345 678.''}
\end{center}

Kết quả phát hiện của module được trình bày trong Bảng~\ref{tab:regex_example}.

\begin{table}[H]
      \centering
      \caption{Kết quả phát hiện thực thể bằng regex detector}
      \label{tab:regex_example}
      \renewcommand{\arraystretch}{1.3}
      \begin{tabular}{l l c c}
            \toprule
            \textbf{Loại}  & \textbf{Text}            & \textbf{start} & \textbf{end} \\
            \midrule
            \texttt{EMAIL} & \texttt{john@a.com}      & 6              & 16           \\
            \texttt{DATE}  & \texttt{01/02/2024}      & 23             & 33           \\
            \texttt{PHONE} & \texttt{+84 912 345 678} & 40             & 55           \\
            \bottomrule
      \end{tabular}
\end{table}

Trong ví dụ trên, ba entity không chồng lấn nên tất cả đều được giữ lại. Nếu
có URL chứa segment giống ngày (ví dụ
\texttt{https://example.com/2024/01/02}), segment đó sẽ bị loại nhờ exclusion
spans và chỉ \texttt{URL} được giữ.
